\section{Piano di documentazione}

Il piano di documentazione definisce il contenuto e la struttura dei documenti del progetto QuantSys Institutional.

\subsection{Documenti accademici (cartella \texttt{PB/})}

Questi documenti sono prodotti in LaTeX e versionati nella cartella \texttt{PB/} della repository:

\begin{itemize}
  \item \textbf{Analisi dei requisiti} (\texttt{analisi\_dei\_requisiti/}): descrive i requisiti funzionali e non funzionali del sistema, i casi d'uso, le interfacce e la tracciabilità;
  \item \textbf{Manuale utente} (\texttt{manuale\_utente/}): guida all'installazione, configurazione, avvio e troubleshooting del sistema;
  \item \textbf{Norme di progetto} (\texttt{norme\_di\_progetto/}): documento corrente. Definisce processi, convenzioni, responsabilità e best practices;
  \item \textbf{Specifica tecnica} (\texttt{specifica\_tecnica/}): descrive l'architettura del sistema, i moduli principali, le interfacce interne e le scelte tecnologiche;
  \item \textbf{Verbali} (\texttt{verbali/}): resoconti degli incontri di team e delle decisioni prese.
\end{itemize}

\subsection{Documenti operativi (radice repository)}

Questi sono documenti Markdown che evolvono con il codice e descrivono aspetti operativi e di sviluppo:

\begin{itemize}
  \item \textbf{\texttt{README.md}}: punto di ingresso per nuovi sviluppatori. Contiene descrizione, quick-start e link agli altri documenti;
  \item \textbf{\texttt{OBIETTIVI.md}}: obiettivi misurabili del sistema (CAGR netto, Sharpe target, alpha vs benchmark, capacity);
  \item \textbf{\texttt{GUIDE\_STRATEGIE\_CUSTOM.md}}: guida tecnica per sviluppatori che vogliono implementare una nuova strategia nel framework plugin. Documento esaustivo che copre struttura cartella, descriptor, indicator pool, universe filtering, exit policy, test e troubleshooting;
  \item \textbf{\texttt{PITCH\_FOF.md}}: documento di pitch verso Fund of Funds, costruito a partire dai numeri OOS validati;
  \item \textbf{\texttt{ROADMAP\_FASE2\_v2.md}}: piano per l'ampliamento dell'ensemble con nuove strategie;
  \item \textbf{\texttt{STRATEGIE\_FX\_CANDIDATE.md}}: proposta di sei strategie FX scrivibili nel framework plugin, da validare in walk-forward.
\end{itemize}

\subsection{Documentazione interna alle strategie}

Ogni strategia plugin (es. \texttt{strategies/momentum/}) può contenere al proprio interno:

\begin{itemize}
  \item \texttt{README.md}: razionale economico, riferimenti letteratura, parametri default, risultati OOS;
  \item File Markdown con prefisso \texttt{\_} (es. \texttt{\_FAILED.md}, \texttt{\_NOTE.md}): documenti interni che il sistema di discovery ignora, utilizzati per registrare tentativi falliti o note di lavoro. La convenzione del prefisso \texttt{\_} è descritta in dettaglio nel Capitolo \ref{chap:strategie}.
\end{itemize}

\subsection{Aggiornamento e versionamento}

Ogni documento accademico deve essere aggiornato regolarmente. Le modifiche ai documenti sono effettuate tramite pull request, esattamente come per il codice, in modo da garantire revisione e tracciabilità.

I documenti Markdown operativi (radice repository) evolvono naturalmente con il codice e vengono aggiornati nelle stesse pull request che modificano le funzionalità correlate.
